% 04-queue-management.tex - Queue Management

\chapter{Queue Management}
\label{sec:queue-management}

This section describes the persistent queue system used by \SC{} to manage copy and delete operations.

\section{Queue Architecture}

\SC{} maintains one logical queue per \DR{} (DR1--DR5), all backed by a single shared SQLite database file (\file{smartcopy.db}). Each queue operates as a FIFO (First In, First Out) structure with the following properties:

\begin{itemize}
    \item \textbf{Persistence}: All queue items are persisted to disk via SQLite, surviving system restarts and crashes
    \item \textbf{Thread Safety}: Database access is protected by a reentrant lock (\code{threading.RLock}) shared across all queue instances
    \item \textbf{WAL Mode}: The database uses Write-Ahead Logging for improved concurrent read/write performance
    \item \textbf{Transactional}: All queue mutations use \code{IMMEDIATE} transactions for atomicity
    \item \textbf{Integrity Recovery}: On startup, any items left in ``processing'' state (from a crash) are automatically reset to ``pending''
\end{itemize}

\section{Database Schema}

The queue database schema is shown in Table~\ref{tab:db-schema}.

\begin{table}[htbp]
    \centering
    \caption{Queue Database Schema (\code{queue\_items} Table)}
    \label{tab:db-schema}
    \begin{tabular}{llp{6cm}}
        \toprule
        \textbf{Column} & \textbf{Type} & \textbf{Description} \\
        \midrule
        \code{id} & INTEGER & Auto-incrementing primary key \\
        \code{queue\_name} & TEXT & Queue identifier (e.g., ``DR1'', ``DR2'') \\
        \code{host} & TEXT & Source hostname \\
        \code{hostpath} & TEXT & Source file path \\
        \code{dest} & TEXT & Destination hostname \\
        \code{destpath} & TEXT & Destination file path (or delete marker) \\
        \code{filesize} & INTEGER & File size in bytes (default: 0) \\
        \code{command\_id} & INTEGER & Unique copy/delete ID \\
        \code{retry\_count} & INTEGER & Number of retry attempts (default: 0) \\
        \code{last\_try} & REAL & Timestamp of last attempt (default: 0.0) \\
        \code{status} & TEXT & Item status (default: ``pending'') \\
        \code{created\_at} & TIMESTAMP & Creation timestamp \\
        \bottomrule
    \end{tabular}
\end{table}

Two indexes are maintained for query performance:
\begin{itemize}
    \item \code{idx\_queue\_status} on \code{(queue\_name, status)} --- used for queue operations and statistics
    \item \code{idx\_created} on \code{(created\_at)} --- used for time-based queries
\end{itemize}

\section{Item Status Lifecycle}

Queue items progress through the statuses shown in Table~\ref{tab:item-status}.

\begin{table}[htbp]
    \centering
    \caption{Queue Item Statuses}
    \label{tab:item-status}
    \begin{tabular}{lp{9cm}}
        \toprule
        \textbf{Status} & \textbf{Description} \\
        \midrule
        \code{pending} & Item is waiting in the queue to be processed \\
        \code{processing} & Item has been dequeued and is currently being copied/deleted \\
        \code{completed} & Transfer completed successfully; source file eligible for purging \\
        \code{failed} & Transfer failed after exhausting all retry attempts or the source file no longer exists \\
        \bottomrule
    \end{tabular}
\end{table}

The lifecycle of a queue item is:

\begin{enumerate}
    \item Item is added via \cmd{SCP} or \cmd{SRM} with status \code{pending}
    \item When dequeued for processing, status changes to \code{processing}
    \item On completion, the \code{processing} row is deleted and a new \code{completed} row is inserted (for successful copies) or a \code{failed} row (for permanent failures)
    \item If a copy fails but retries remain, the item is re-added as \code{pending} with an incremented \code{retry\_count}
\end{enumerate}

\section{Retry Logic}

Failed copy operations are retried according to the following rules:

\begin{itemize}
    \item The maximum number of retries is controlled by the \code{max\_retry} configuration parameter (default: 7)
    \item The minimum wait time between retries is controlled by the \code{wait\_retry} configuration parameter (default: 24 hours)
    \item If a retry is attempted too soon (within \code{wait\_retry} hours of the last attempt), the item is returned to the end of the queue
    \item A copy is permanently failed if:
    \begin{itemize}
        \item The source file no longer exists, or
        \item The retry count has reached \code{max\_retry}
    \end{itemize}
    \item Delete operations (via \cmd{SRM}) have their retry count incremented by an additional 100 (on top of the normal increment of 1) to distinguish them from copy retries
\end{itemize}

\section{Remote Copy Exclusivity}
\label{sec:remote-exclusivity}

To prevent network saturation, only one remote copy may be active at any time across all \DR{} queues. In this context, a copy is considered remote when the source and destination hostnames differ (regardless of whether the destination is on the local network or behind a gateway). This is a simpler test than the gateway detection used for bandwidth limiting (Section~\ref{sec:bw-gateway-detection}). Exclusivity is enforced by a shared \code{threading.Semaphore}. When a \DR{} queue manager encounters a remote copy item:

\begin{enumerate}
    \item It attempts to acquire the remote copy lock (non-blocking)
    \item If the lock is unavailable, the item is returned to the end of the queue and processing is deferred
    \item If the lock is acquired, the remote copy proceeds with the configured bandwidth limit
    \item The lock is released when the copy completes (successfully or with failure)
\end{enumerate}

\section{Delete Markers}

Delete operations use special destination path markers to distinguish them from copy operations. When an \cmd{SRM} command is received, the destination path is set to one of two marker values:

\begin{description}
    \item[\code{smartcopy\_queue\_delete\_this\_file}] Deferred delete: the file is marked as complete immediately and will be deleted during the next purge cycle.
    \item[\code{smartcopy\_now\_delete\_this\_file}] Immediate delete (via \code{-tNOW}): the file is deleted via \code{rm -f} as soon as the item is processed.
\end{description}

\section{Purge Operations}

The purge process runs daily at approximately 18:00~UTC and operates as follows:

\begin{enumerate}
    \item Collect all items with status \code{completed} for the \DR{}
    \item Sum the total file sizes
    \item If the total exceeds the \code{purge\_size} threshold (default: 1~TB), proceed with deletion
    \item For each completed file, execute \code{sudo rm -f} on the \DR{} via SSH
    \item Files that could not be deleted (e.g., due to the queue being paused) are retained for the next purge cycle
    \item After purging, purge all \code{completed} records from the database
\end{enumerate}

\begin{calloutBox}{Implementation Note}{blue}
Spectrometer files (paths containing \file{DROS/Spec}) are only added to the completed list after being successfully copied to the remote archive. This prevents data loss if a spectrometer file is copied locally but not yet archived remotely.
\end{calloutBox}
